\slajdid{03-s010b}
\begin{frame}[label=03-s010b]{Matrična i piramidalna struktura: poređenje}
% element: 03-s010b:e03
Za dekoder \(6/64\) realizovan dekoderima 2/4:
\par\vspace{2mm}
\begin{columns}[T,onlytextwidth]
\begin{column}{.475\textwidth}
\textbf{Matrična struktura}\par\vspace{1mm}
16 izlaznih dekodera + dekoder vrste + dekoder kolone
\[\textcolor{DEisticanje}{16+1+1=18\text{ komponenti}}\]
\end{column}
\begin{column}{.475\textwidth}
\textbf{Piramidalna struktura}\par\vspace{1mm}
16 dekodera u izlaznom + 4 u srednjem + 1 u prvom nivou
\[16+4+1=21\text{ komponenta}\]
\end{column}
\end{columns}
\par\vspace{2mm}
% element: 03-s010b:e04
\Rezultat{\textbf{Matrična struktura ima dva nivoa, pa i manje kašnjenje.}
U piramidalnoj strukturi broj nivoa zavisi od veličine dekodera;
za ovaj primer potrebna su tri nivoa.}
\par\vspace{2mm}
Pri većem broju izlaza prednost matrične strukture u broju komponenti postaje značajnija,
ali zahteva \textbf{specifičnije komponente sa više CS ulaza}.
Moguće je kombinovati matričnu i piramidalnu strukturu.
\end{frame}
